\documentclass[11pt]{article}

\usepackage[margin=1in]{geometry}
\usepackage{amsmath,amssymb}
\usepackage{booktabs}
\usepackage{hyperref}
\usepackage{enumitem}
\usepackage{xcolor}
\usepackage{cite}
\usepackage{url}

\hypersetup{
    colorlinks=true,
    linkcolor=blue,
    citecolor=blue,
    urlcolor=blue
}

\setlength{\parindent}{0pt}
\setlength{\parskip}{0.6em}

\title{Paper Review of Phase 1 Report:\\Lightweight Spike-Based DOA Sector Inference for Embedded Edge AI}
\author{Hossein Molhem\\Kennesaw State University\\EE 6900 Special Topics: Embedded AI for Autonomous Robots}
\date{\today}

\begin{document}
\maketitle

\begin{abstract}
This paper review summarizes and evaluates the Phase 1 report on lightweight spike-based direction-of-arrival (DOA) sector inference for embedded edge AI. The reviewed work presents a reproducible Python prototype that connects a physics-based uniform linear array (ULA) signal model, sample covariance matrix estimation, beamspace feature extraction, nearest-centroid baseline classification, rate-based spike encoding, and leaky-integrate-and-fire (LIF)-inspired temporal inference. The review assesses the technical motivation, methodology, experimental design, quantitative results, engineering interpretation, limitations, and recommended future work. The main conclusion is that Phase 1 provides a sound baseline framework for hardware-aware embedded-AI exploration. The spike-based path does not improve raw classification accuracy over the continuous nearest-centroid baseline; however, it introduces sparse temporal inference, measurable stable-decision latency, and tunable accuracy--latency behavior, which are directly relevant to embedded and neuromorphic edge-inference studies.
\end{abstract}

\section{Review Scope}

This review is based on the Phase 1 project report titled \emph{DOA-Sector Inference with LIF-Inspired Edge Processing}. The report documents a staged embedded-AI prototype for sector-level DOA classification. The work is not presented as a new high-resolution DOA estimator. Instead, it is positioned as an engineering baseline that links array signal processing to spike-based temporal inference and later hardware-aware deployment.

The reviewed report is technically organized around three layers: array signal processing, classical machine-learning baseline classification, and sparse temporal inference. This structure is appropriate because it separates the sensing model, feature representation, and edge-inference mechanism. The result is a controlled framework that can be studied quantitatively before attempting FPGA, PYNQ, or Vitis HLS implementation.

\section{Summary of the Reviewed Work}

The Phase 1 prototype generates synthetic narrowband ULA data for a single source over a controlled angular and SNR grid. For each sample, the received snapshot matrix is modeled using a ULA steering vector, a complex narrowband source waveform, and additive complex Gaussian noise. The sample covariance matrix is computed as
\begin{equation}
    R_{xx}=\frac{1}{M}XX^{H},
\end{equation}
where $M=128$ snapshots are used per sample. Beamspace energy features are extracted from the covariance matrix over a fixed angular grid and then normalized using training-set statistics.

The classification task is intentionally framed as sector-level inference rather than continuous-angle estimation. A nearest-centroid classifier is used as the deterministic non-spiking baseline. The normalized beamspace features are then converted into binary spike trains through rate-based encoding. Finally, a LIF-inspired classifier accumulates spike-driven evidence over time according to
\begin{equation}
    v_c[t]=\alpha v_c[t-1]+i_c[t],
\end{equation}
where $v_c[t]$ is the class membrane state, $\alpha$ is the leak factor, and $i_c[t]$ is the spike-driven input current. Stable-decision latency is measured as the earliest time step after which the final decision remains unchanged.

\section{Technical Motivation}

The motivation is strong and appropriate for an embedded-AI course. Edge systems are often constrained by compute resources, energy budget, memory footprint, and latency. A conventional high-capacity neural network may not be the best first choice for such platforms, especially when the sensing problem has an interpretable physics-based structure. The selected DOA-sector task is well chosen because it naturally links electromagnetic sensing, array signal processing, compact feature extraction, and hardware-aware inference.

The report also benefits from a clear connection to prior author work on FPGA/HLS sample covariance matrix acceleration for DOA preprocessing~\cite{molhem2026scmhls}. This connection strengthens the engineering narrative: the current project is not an isolated Python experiment, but part of a broader path from covariance-based DOA preprocessing toward embedded and FPGA-aware inference systems.

\section{Methodological Review}

The methodology is defensible for a Phase 1 baseline. The use of a ULA model creates a physically meaningful data source. The sample covariance matrix captures spatial correlation information. Beamspace features reduce the covariance structure into a compact angular representation, which is a reasonable design choice for coarse sector classification. The nearest-centroid classifier provides a simple and interpretable baseline before introducing stochastic spike encoding.

The LIF-inspired inference model is also appropriately scoped. The report does not overclaim biological realism. Instead, the LIF path is framed as a lightweight hardware-facing approximation: each class integrates spike-driven evidence over time with leakage. This is a useful abstraction because it creates a temporal decision process and allows latency to become an explicit design metric.

A key methodological strength is the explicit comparison between a continuous-valued baseline and a spike-based path. Without this comparison, it would be difficult to know whether the spike encoding preserves useful information. The report correctly treats the nearest-centroid baseline as the accuracy reference point.

\section{Experimental Design}

The dataset configuration is compact but sufficient for Phase 1. It uses 8 array elements, 128 snapshots per sample, half-wavelength spacing, one source, four DOA sectors, seven SNR values from $-10$ dB to $20$ dB, 2800 total samples, and a stratified 70/30 train/test split. This provides enough variation to evaluate accuracy versus SNR while keeping the experiment reproducible and computationally lightweight.

The evaluation metrics are appropriate for the stated objectives. The report includes overall accuracy, accuracy versus SNR, spike-count statistics, stable-decision latency, and an accuracy--latency trade-off study. This is important because embedded inference should not be judged only by classification accuracy. A lower-accuracy model may still be valuable if it exposes latency control, sparsity, and implementation-friendly temporal structure.

\section{Results Review}

The nearest-centroid baseline achieved an overall test accuracy of $0.869048$, or approximately $86.90\%$. This result demonstrates that the compact beamspace features contain meaningful information for coarse DOA-sector inference.

The fixed LIF-inspired configuration used $T=20$ and leak $=0.97$. It achieved an accuracy of $0.825000$, or $82.50\%$, with a mean stable-decision latency of $6.423810$ time steps. This confirms that the spike-based path preserves much of the baseline information, but with a measurable reduction in raw accuracy.

The Run 005 parameter sweep evaluated $T\in\{10,20,30,40\}$ and leak $\in\{0.90,0.95,0.97,0.99\}$, producing 16 configurations. The best LIF accuracy was $0.841667$ at $T=40$ and leak $=0.99$, with mean latency $8.396429$ time steps. The lowest-latency configuration occurred at $T=10$ and leak $=0.99$, with accuracy $0.776190$ and mean latency $3.782143$ time steps.

\begin{table}[h]
\centering
\caption{Key Phase 1 Results}
\begin{tabular}{lccc}
\toprule
Case & Time window & Leak & Accuracy \\
\midrule
Nearest-centroid baseline & -- & -- & 0.869048 \\
Fixed LIF configuration & 20 & 0.97 & 0.825000 \\
Best LIF sweep configuration & 40 & 0.99 & 0.841667 \\
Lowest-latency LIF configuration & 10 & 0.99 & 0.776190 \\
\bottomrule
\end{tabular}
\end{table}

\section{Engineering Interpretation}

The most important interpretation is that the spike-based path is not an accuracy-improving replacement for the continuous nearest-centroid baseline. The baseline remains the best raw-accuracy classifier. This is not a failure of the project; rather, it is the central engineering result. Rate-based spike encoding introduces stochastic quantization, and temporal accumulation creates a design space in which accuracy, latency, and sparsity must be balanced.

The value of the spike-based formulation is architectural. It converts continuous beamspace features into sparse temporal events. It exposes stable-decision latency as a measurable output. It allows the designer to adjust the spike time window and leak factor to trade accuracy against temporal budget. These properties are directly aligned with embedded-AI and neuromorphic-inference motivations.

\section{Limitations}

The Phase 1 study is intentionally controlled, and several limitations should be noted. First, the dataset uses a single-source narrowband ULA model, which is simpler than many real radar or array-processing scenarios. Second, the task is sector classification rather than high-resolution DOA estimation. Third, the spike encoder is rate-based and stochastic, so performance depends on time-window length and random spike realization. Fourth, the LIF model is a simplified class-level accumulator rather than a detailed spiking neural network. Finally, the current implementation is a Python reference prototype, not yet a fixed-point or deployed hardware implementation.

These limitations are acceptable for Phase 1 because the purpose is not to solve the full radar/DOA deployment problem. The purpose is to establish a clean baseline and identify a meaningful hardware-aware direction.

\section{Recommended Future Work}

The most important next step is hardware-aware refactoring. The covariance, beamspace, spike encoding, and temporal accumulation stages should be evaluated for fixed-point representation and implementation cost. The prior SCM HLS accelerator provides a natural starting point for the covariance stage. Later phases can integrate the Python DOA/SNN workflow with a fixed-point HLS or PYNQ-oriented backend.

Additional recommended work includes multi-source DOA cases, randomized waveform and noise studies, richer SNR and angular distributions, confusion-matrix analysis by sector, energy or operation-count estimates, and systematic fixed-point quantization experiments. A smaller FPGA demonstrator could also be developed to show the complete path from signal-processing theory to hardware-aware edge inference.

\section{Overall Assessment}

Overall, the Phase 1 report is technically sound, well scoped, and defensible as an embedded-AI project. Its strongest feature is the integration of array signal processing, interpretable baseline classification, spike-based temporal inference, and latency-aware evaluation. The results are not overstated: the report clearly acknowledges that the spike-based path does not outperform the continuous baseline in raw accuracy. Instead, it correctly frames the contribution as an architectural bridge toward sparse, temporal, hardware-aware inference.

For presentation to Dr. Feng, the central message should be concise: Phase 1 establishes a reproducible DOA-sector inference baseline; the continuous beamspace classifier provides strong accuracy; the spike-based path introduces a controlled accuracy--latency trade-off; and the project now has a clear path toward FPGA/HLS-aware embedded implementation.

\clearpage
\bibliographystyle{IEEEtran}
\bibliography{reference}

\end{document}
